\section{Stabilité et propriétés discrètes démontrées}
\label{sec:discrete-properties}

\subsection{Condition suffisante et coefficients de mise à jour}

Écrivons $\mathcal N(P)$ pour l'ensemble des voisins d'un pixel $P$.
Si l'interface $PQ$ est horizontale, posons $h_{PQ}=\dx$ ; si elle
est verticale, $h_{PQ}=\dy$. La conductance partagée est
$a^n_{PQ}=a^n_{QP}$, égale à un pour la chaleur et comprise entre zéro
et un pour la variante Perona--Malik. Le schéma peut s'écrire
\begin{equation}
 U^{n+1}_P=
 \left(1-\sum_{Q\in\mathcal N(P)}\alpha^n_{PQ}\right)U^n_P
 +\sum_{Q\in\mathcal N(P)}\alpha^n_{PQ}U^n_Q,
 \qquad
 \alpha^n_{PQ}=\frac{\dt\,a^n_{PQ}}{h_{PQ}^2}.
 \label{eq:convex-update}
\end{equation}
Tous les $\alpha^n_{PQ}$ sont non négatifs et
\begin{equation}
 \sum_{Q\in\mathcal N(P)}\alpha^n_{PQ}
 \leq 2\dt\left(\frac1{\dx^2}+\frac1{\dy^2}\right).
 \label{eq:weight-bound}
\end{equation}
La présence de moins de voisins sur le bord ne peut qu'abaisser cette
somme. La condition suffisante
\begin{equation}
 \boxed{\dt\left(\frac1{\dx^2}+\frac1{\dy^2}\right)\leq\frac12}
 \label{eq:cfl}
\end{equation}
rend donc le coefficient central non négatif. Les coefficients de
\eqref{eq:convex-update} somment à un. Avec $\dx=\dy=1$ et
$\dt=0{,}20$, le membre de gauche vaut $0{,}40<0{,}50$ ; la somme
maximale des poids des voisins vaut $0{,}80$ et le poids central
est au moins $0{,}20$.

Cette condition est volontairement suffisante et commune aux trois
solveurs ; des conductances faibles pourraient autoriser ponctuellement
une borne moins restrictive, mais aucun pas adaptatif n'est utilisé.
Le programme rejette un pas invalide au lieu de le réduire silencieusement.
La comparaison de la borne inclut seulement une marge de
$32\varepsilon_{\mathrm{machine}}$ pour l'arrondi flottant du calcul
de stabilité, et non une permission de dépasser matériellement la
condition.

\subsection{Conservation des constantes et de la moyenne}

\paragraph{Constantes.}
Si $U^n_P=b$ pour tous les pixels, chaque différence d'interface est
nulle. Toutes les contributions $Q^x,Q^y$ sont donc nulles, quelle
que soit la conductance. Ainsi $U^{n+1}_P=b$. Par récurrence, une image
constante demeure constante pour tout nombre d'étapes. Cette propriété
n'impose pas que $b$ appartienne à $[0,1]$.

\paragraph{Somme.}
En sommant \eqref{eq:explicit-euler} sur les pixels, une interface
intérieure contribue une fois par $+\dt Q$ et une fois par
$-\dt Q$. Les deux termes s'annulent. Il n'existe aucune contribution
extérieure selon \eqref{eq:boundary-edges}. En arithmétique exacte,
\begin{equation}
 \sum_P U^{n+1}_P=\sum_P U^n_P=\sum_P f_P,
 \qquad \overline U^{n+1}=\overline U^n=\overline f.
 \label{eq:discrete-mass}
\end{equation}
Sur une grille uniforme, la masse pondérée par l'aire des cellules,
$\dx\dy\sum_P U_P$, est également conservée. Le résultat reste
valable lorsque les conductances dépendent de l'état : l'annulation
requiert leur partage sur l'interface, pas leur constance dans le
temps. En machine, les additions dans des cases distinctes puis la
sommation peuvent introduire un arrondi ; les tests contrôlent donc
la conservation à une tolérance absolue adaptée.

\subsection{Principe du maximum discret}

Sous \eqref{eq:cfl}, chaque valeur nouvelle est une combinaison
convexe de valeurs anciennes. Si $m_n=\min_P U^n_P$ et
$M_n=\max_P U^n_P$, alors
\begin{equation}
 m_n\leq U^{n+1}_P\leq M_n\quad\text{pour tout }P,
 \qquad m_{n+1}\geq m_n,\quad M_{n+1}\leq M_n.
 \label{eq:discrete-maximum}
\end{equation}
Par récurrence, $\min f\leq U^n_P\leq\max f$. La démonstration vaut
pour Perona--Malik malgré sa non-linéarité : à l'étape considérée,
les coefficients calculés depuis $U^n$ vérifient les mêmes bornes.

Il ne s'agit pas d'une borne dans $[0,1]$. Si le bruit place une valeur
initiale à $-0{,}1$ ou à $1{,}1$, le principe du maximum concerne cet
intervalle initial. Il est donc compatible avec l'absence de troncature.
Il ne prouve pas non plus une contraction entre deux solutions non
linéaires distinctes : leurs conductances dépendent de leurs états
respectifs. Aucune telle contraction n'est déduite de
\eqref{eq:convex-update}.

\subsection{Dissipation quadratique de la chaleur sur la grille}

Pour la chaleur seule, $L_h=L$ est une matrice réelle symétrique.
En regroupant les contributions par interface, on obtient
\begin{equation}
 \langle V,LV\rangle
 =-\sum_{\{P,Q\}\in\mathcal E}
    \frac{(V_Q-V_P)^2}{h_{PQ}^2}\leq0,
 \label{eq:linear-discrete-form}
\end{equation}
où $\mathcal E$ contient chaque interface intérieure une fois.
La matrice est négative semi-définie et annule les constantes.
Pour le rectangle à flux nul, ses valeurs propres sont comprises
entre $-4/\dx^2-4/\dy^2$ et zéro, comme le confirme la formule
explicite des modes de la section~\ref{sec:verification}.

Posons $V^n=U^n-\overline f$. Dans une base orthonormée de vecteurs
propres, chaque composante d'Euler est multipliée par
$1+\dt\lambda$, avec $\lambda\leq0$. Sous \eqref{eq:cfl},
$-1\leq1+\dt\lambda\leq1$. Il en résulte
\begin{equation}
 \sum_P (U^{n+1}_P-\overline f)^2
 \leq\sum_P (U^n_P-\overline f)^2.
 \label{eq:linear-discrete-l2}
\end{equation}
Cette dissipation est celle d'une quantité discrète précisément
définie, et non une affirmation sur l'erreur par rapport à l'image
propre. Le test existant emploie un état non constant et vérifie une
diminution stricte après plusieurs pas. La preuve spectrale précédente
est linéaire ; elle n'est pas transférée ici à une énergie non linéaire
de Perona--Malik.

\subsection{Ce que ces propriétés ne démontrent pas}

La \emph{stabilité discrète} examinée ici est le maintien de bornes
et de propriétés conservatrices à grille fixée. La \emph{consistance}
compare un opérateur discret à un opérateur différentiel sur des
fonctions régulières. La \emph{convergence} demande un passage à la
limite en raffinant la grille et le temps dans un cadre précisé.
Enfin, la \emph{bonne position} du problème continu concerne
l'existence, l'unicité et la dépendance vis-à-vis des données.
Les preuves précédentes ne rendent pas ces quatre notions
interchangeables. En particulier, elles ne prouvent pas une convergence
de la variante directionnelle vers le modèle scalaire
\eqref{eq:pm-continuous}, dont elle n'approche pas le même opérateur
dans sa limite formelle intérieure.
